% GENERATED FILE. Edit canonical lessons, then run npm run generate:books.
% canonical-source-hash: fnv1a64:b57f9f29112c179f
% canonical-lessons: FR-C363-manifestation, FR-C363-manifester, FR-C363-manifestant, FR-C363-banderole, FR-C363-protester

\chapter{Dans la rue}
\label{ch:fr-dans-la-rue}
\hlchaptermodality{363}

\begin{chapteropening}
\textbf{By the end of this chapter:} \emph{I can describe a protest: a demonstration, demonstrators, a banner, and people who demonstrate and protest.}

Complete the last lesson of chapter 363: I can describe a protest: a demonstration, demonstrators, a banner, and people who demonstrate and protest.
\end{chapteropening}

\section[la manifestation]{la manifestation — the demonstration, the protest march; the event}

\label{lesson:FR-C363-manifestation}



\begin{quote}
Before the new one: say the French for the freedom, then the French for the equality.
\end{quote}

\subsection*{You'll want to know: la manifestation}
\textbf{la manifestation} — \textquotedblleft{}the demonstration, the protest march; the event\textquotedblright{}.

\begin{quote}
Il y a une manifestation dans le centre-ville samedi.
\end{quote}

\begin{quote}
\emph{There is a demonstration in the town centre on Saturday.}
\end{quote}

Often shortened to \emph{la manif}. Also any public event: \emph{une manifestation sportive}.

\begin{cousinweb}
\emph{manifestation} comes from Late Latin \emph{manifēstātiō} \textquotedblleft{}a making plain\textquotedblright{} (EtymDB), from \emph{manifestus} \textquotedblleft{}clear, plain\textquotedblright{} (Lewis \& Short).
\end{cousinweb}

\begin{grammarlens}[title={What you've built}]
One more word for telling, explaining and describing at B1.
\end{grammarlens}

\subsection*{Your turn}
\begin{itemize}
\raggedright
  \item \emph{Say it:} \emph{la manifestation}
  \item \emph{Say it:} \emph{la manifestation}, once more
  \item \emph{Say it:} a sentence of your own with \emph{la manifestation}
  \item \emph{Recall:} say the French for the freedom, then the French for the equality, then say \emph{la manifestation} again
\end{itemize}

\begin{tcolorbox}[breakable,colback=teal!4,colframe=teal!35!black,title={Before you move on}]
What does la manifestation mean? (\textquotedblleft{}The demonstration, the protest march; the event\textquotedblright{}.) Say it once more.
\end{tcolorbox}
\section[manifester]{manifester — to demonstrate (in the street); to show}

\label{lesson:FR-C363-manifester}



\begin{quote}
Before the new one: say the French for the equality, then the French for the demonstration.
\end{quote}

\subsection*{You'll want to know: manifester}
\textbf{manifester} — \textquotedblleft{}to demonstrate (in the street); to show\textquotedblright{}.

\begin{quote}
Les étudiants manifestent contre la nouvelle loi.
\end{quote}

\begin{quote}
\emph{The students are demonstrating against the new law.}
\end{quote}

\emph{Manifester sa joie} — to show your joy.

\begin{cousinweb}
\emph{manifester} is a learned borrowing from Latin \emph{manifēstāre} \textquotedblleft{}to make plain, to show\textquotedblright{} (Lewis \& Short), from \emph{manifestus} \textquotedblleft{}clear, plain\textquotedblright{}.
\end{cousinweb}

\begin{grammarlens}[title={What you've built}]
One more word for telling, explaining and describing at B1.
\end{grammarlens}

\subsection*{Your turn}
\begin{itemize}
\raggedright
  \item \emph{Say it:} \emph{manifester}
  \item \emph{Say it:} \emph{manifester}, once more
  \item \emph{Say it:} a sentence of your own with \emph{manifester}
  \item \emph{Recall:} say the French for the equality, then the French for the demonstration, then say \emph{manifester} again
\end{itemize}

\begin{tcolorbox}[breakable,colback=teal!4,colframe=teal!35!black,title={Before you move on}]
What does manifester mean? (\textquotedblleft{}To demonstrate (in the street); to show\textquotedblright{}.) Say it once more.
\end{tcolorbox}
\section[le manifestant]{le manifestant — the demonstrator, the marcher}

\label{lesson:FR-C363-manifestant}



\begin{quote}
Before the new one: say the French for the demonstration, then the French for to demonstrate.
\end{quote}

\subsection*{You'll want to know: le manifestant}
\textbf{le manifestant} — \textquotedblleft{}the demonstrator, the marcher\textquotedblright{}.

\begin{quote}
Les manifestants marchent vers la mairie.
\end{quote}

\begin{quote}
\emph{The demonstrators are marching towards the town hall.}
\end{quote}

Feminine \emph{manifestante}.

\begin{cousinweb}
\emph{manifestant} is the present participle of \emph{manifester} used as a noun. The verb comes from Latin \emph{manifēstāre} \textquotedblleft{}to make plain\textquotedblright{}, from \emph{manifestus} \textquotedblleft{}clear\textquotedblright{} (Lewis \& Short).
\end{cousinweb}

\begin{grammarlens}[title={What you've built}]
One more word for telling, explaining and describing at B1.
\end{grammarlens}

\subsection*{Your turn}
\begin{itemize}
\raggedright
  \item \emph{Say it:} \emph{le manifestant}
  \item \emph{Say it:} \emph{le manifestant}, once more
  \item \emph{Say it:} a sentence of your own with \emph{le manifestant}
  \item \emph{Recall:} say the French for the demonstration, then the French for to demonstrate, then say \emph{le manifestant} again
\end{itemize}

\begin{tcolorbox}[breakable,colback=teal!4,colframe=teal!35!black,title={Before you move on}]
What does le manifestant mean? (\textquotedblleft{}The demonstrator, the marcher\textquotedblright{}.) Say it once more.
\end{tcolorbox}
\section[la banderole]{la banderole — the banner}

\label{lesson:FR-C363-banderole}



\begin{quote}
Before the new one: say the French for to demonstrate, then the French for the demonstrator.
\end{quote}

\subsection*{You'll want to know: la banderole}
\textbf{la banderole} — \textquotedblleft{}the banner\textquotedblright{}.

\begin{quote}
Les manifestants portent une grande banderole.
\end{quote}

\begin{quote}
\emph{The demonstrators are carrying a big banner.}
\end{quote}

A long strip of cloth with words on it.

\begin{cousinweb}
\emph{banderole} was borrowed from Italian \emph{banderuola} \textquotedblleft{}a little flag\textquotedblright{} (EtymDB), a small form of \emph{bandiera} \textquotedblleft{}flag\textquotedblright{}.
\end{cousinweb}

\begin{grammarlens}[title={What you've built}]
One more word for telling, explaining and describing at B1.
\end{grammarlens}

\subsection*{Your turn}
\begin{itemize}
\raggedright
  \item \emph{Say it:} \emph{la banderole}
  \item \emph{Say it:} \emph{la banderole}, once more
  \item \emph{Say it:} a sentence of your own with \emph{la banderole}
  \item \emph{Recall:} say the French for to demonstrate, then the French for the demonstrator, then say \emph{la banderole} again
\end{itemize}

\begin{tcolorbox}[breakable,colback=teal!4,colframe=teal!35!black,title={Before you move on}]
What does la banderole mean? (\textquotedblleft{}The banner\textquotedblright{}.) Say it once more.
\end{tcolorbox}
\section[protester]{protester — to protest}

\label{lesson:FR-C363-protester}



\begin{quote}
Before the new one: say the French for the demonstrator, then the French for the banner.
\end{quote}

\subsection*{You'll want to know: protester}
\textbf{protester} — \textquotedblleft{}to protest\textquotedblright{}.

\begin{quote}
Les voisins protestent contre le bruit.
\end{quote}

\begin{quote}
\emph{The neighbours are protesting about the noise.}
\end{quote}

\emph{Protester contre quelque chose}.

\begin{cousinweb}
\emph{protester} comes from Latin \emph{prōtestārī} \textquotedblleft{}to declare publicly, to bear witness\textquotedblright{} (EtymDB; Lewis \& Short), built on \emph{testis} \textquotedblleft{}witness\textquotedblright{}.
\end{cousinweb}

\begin{grammarlens}[title={What you've built}]
One more word for telling, explaining and describing at B1.
\end{grammarlens}

\subsection*{Your turn}
\begin{itemize}
\raggedright
  \item \emph{Say it:} \emph{protester}
  \item \emph{Say it:} \emph{protester}, once more
  \item \emph{Say it:} a sentence of your own with \emph{protester}
  \item \emph{Recall:} say the French for the demonstrator, then the French for the banner, then say \emph{protester} again
\end{itemize}

\begin{tcolorbox}[breakable,colback=teal!4,colframe=teal!35!black,title={Before you move on}]
What does protester mean? (\textquotedblleft{}To protest\textquotedblright{}.) Say it once more.
\end{tcolorbox}
